\documentclass[11pt]{article}
\usepackage[a4paper,margin=1in]{geometry}
\usepackage{fontspec}
\usepackage[boldfont=false]{xeCJK}
\setCJKmainfont{FandolSong-Regular.otf}
\usepackage{amsmath}
\usepackage{amssymb}
\usepackage{array}
\usepackage{booktabs}
\usepackage{graphicx}
\usepackage{adjustbox}
\usepackage{enumitem}
\setlist[itemize]{topsep=2pt,partopsep=0pt,parsep=0pt,itemsep=2pt,leftmargin=1.4em}
\setlist[enumerate]{topsep=2pt,partopsep=0pt,parsep=0pt,itemsep=2pt,leftmargin=1.6em}
\usepackage{parskip}
\usepackage{fvextra}
\fvset{breaklines=true,breakanywhere=true,fontsize=\small}
\setlength{\parindent}{0pt}

\begin{document}

\begin{center}
  {\LARGE\bfseries Electrochemistry}\\[0.35em]
  {\large A Level Chemistry}
\end{center}
\vspace{0.6em}

\section*{Electrolysis}

\textbf{Electrolysis} 电解 uses electricity to break down a molten or aqueous \textbf{electrolyte} 电解质. Positive ions move to the negative electrode and negative ions move to the positive electrode.

At each \textbf{electrode} 电极 a half-reaction happens:

\begin{itemize}
\item at the \textbf{cathode} 阴极 (negative): positive ions gain electrons (reduction).

\item at the \textbf{anode} 阳极 (positive): negative ions lose electrons (oxidation).

\end{itemize}

\subsection*{Predicting the products}

\begin{itemize}
\item a \textbf{molten} electrolyte gives the metal at the cathode and the non-metal at the anode.

\item an \textbf{aqueous} electrolyte also contains water. At the cathode, a less reactive metal is released, but for a reactive metal you get hydrogen instead. At the anode you usually get oxygen, but a concentrated halide solution gives the halogen.

\end{itemize}

\subsection*{Calculations}

The charge on one mole of electrons is the \textbf{Faraday constant} 法拉第常量, linked to the \textbf{Avogadro constant} 阿伏伽德罗常量 ($L$) and the charge on one electron ($e$) by $F = Le$.

The charge passed is $Q = It$ (current $\times$ time). Then:

\begin{enumerate}
\item moles of electrons $= Q / F$.

\item use the half-equation to find moles of product.

\item find the mass ($\times M_r$) or the gas volume.

\end{enumerate}

Measuring the mass deposited for a known charge lets you work back to a value of the Avogadro constant.

\section*{Standard electrode potentials and cell potentials}

The \textbf{electrode potential} 电极电势 ($E$) of a half-cell shows how easily it is reduced. We measure it against a reference, under standard conditions, to get the \textbf{standard electrode potential} 标准电极电势 ($E^{\ominus}$), always written as a reduction.

The reference is the \textbf{standard hydrogen electrode} 标准氢电极: hydrogen gas at $1\ \text{atm}$ over platinum in $1\ \text{mol dm}^{-3}$ $\text{H}^+$, defined as exactly $0.00\ \text{V}$.

To measure an $E^{\ominus}$, connect the half-cell to the standard hydrogen electrode and read the voltage. A metal sits in a solution of its ions; for two ions of the same element (such as $\text{Fe}^{3+}/\text{Fe}^{2+}$), a platinum electrode dips into a solution containing both.

\subsection*{Combining half-cells}

The \textbf{standard cell potential} 标准电池电势 is the difference between the two standard electrode potentials:

\[
E^{\ominus}_{\text{cell}} = E^{\ominus}(\text{more positive}) - E^{\ominus}(\text{less positive})
\]

From $E^{\ominus}$ values you can:

\begin{itemize}
\item find the \textbf{polarity}: the more negative electrode is the negative terminal, and electrons flow from it through the external circuit to the positive electrode.

\item judge \textbf{reactivity}: a more positive $E^{\ominus}$ means a better \textbf{oxidising agent} 氧化剂 (easily reduced); a more negative $E^{\ominus}$ means a better \textbf{reducing agent} 还原剂.

\end{itemize}

\subsection*{Feasibility and redox equations}

A reaction is feasible when $E^{\ominus}_{\text{cell}}$ is \textbf{positive}. To build the full equation, take the two \textbf{half-equations} 半反应方程式, reverse the one that is oxidised, and add them so the electrons cancel. This $E^{\ominus}$ test tells you the \textbf{feasibility} 可行性 of the reaction.

You can also link it to free energy: $\Delta G^{\ominus} = -n E^{\ominus}_{\text{cell}} F$, where $n$ is the moles of electrons.

\subsection*{The Nernst equation}

If the concentrations are not standard, the electrode potential changes. Raising the concentration of the \textbf{oxidised} species makes $E$ more positive. The \textbf{Nernst equation} 能斯特方程 gives the exact value:

\[
E = E^{\ominus} + \frac{0.059}{z} \log \frac{[\text{oxidised species}]}{[\text{reduced species}]}
\]

where $z$ is the number of electrons in the half-equation.


\end{document}
